\subsection{Non-Commutativity as a Source of Mass}
  \label{sec:noncommutativity-mass}

  Within the spectral framework established in Chapter~\ref{sec:spectral-admissibility},
  torsion acts as a dynamical constraint that competes with
  relaxation transport through the projection fiber.

  \subsubsection*{Inhibition of Projective Ordering}
    \label{sec:inhibition-relaxation}

    Let $\Omega_r$ denote the effective torsion operator associated
    with the abstract torsion sector index~$r$.
    For the fundamental sector ($r=1$),
    torsion and projective ordering remain spectrally compatible:
    \begin{equation}
      [\Delta^{(0)}_G,\Omega_{1}] = 0 .
    \end{equation}

    For higher torsion sectors ($r\ge 2$), torsional constraints
    become spectrally frustrated:
    \begin{equation}
      [\Delta^{(0)}_G,\Omega_{r}] \neq 0
      \qquad (r\ge 2).
    \end{equation}

    This non-commutativity inhibits uniform projective ordering across the
    projection fiber and induces an effective spectral compression,
    which manifests as increased inertial mass.

    The torsional contribution can be expressed through the spectral
    invariant
    \begin{equation}
      \mathcal{A}(r) \equiv
      \frac12 \ln\!\left(
                     \frac{\det_\zeta(\Delta^{(0)}_G+\Omega_{r})}
                     {\det_\zeta(\Delta^{(0)}_G)}
      \right),
      \label{eq:Aw-fredholm}
    \end{equation}
    where $\det_\zeta$ denotes zeta-regularized determinants.

  \subsubsection*{Dynamical Stability of Torsional Frequencies}
    \label{sec:pisano-ratio}

    For $r=2$, the projection fiber becomes dynamically anisotropic
    and splits into competing sectors
    \[
      \Pi=\Pi_{\parallel}\oplus\Pi_{\perp}.
    \]

    The quantities $\lambda_{\parallel}$ and $\lambda_{\perp}$ are
    effective torsional rotation frequencies generated by $\Omega_2$
    within the first multiplet.
    They are extracted from the commutator flow on observables
    restricted to $\mathcal{H}_{k=1}$.

    Dynamical stability favors frequency ratios that remain robust
    under perturbations.
    In analogy with KAM-type criteria, the most stable admissible
    ratio corresponds to the golden ratio
    \begin{equation}
      \frac{\lambda_{\parallel}}{\lambda_{\perp}}
      = \varphi ,
      \qquad
      \varphi=\frac{1+\sqrt5}{2}.
      \label{eq:beta-golden}
    \end{equation}

    The appearance of $\varphi$ therefore reflects a stability
    selection rule rather than a static spectral extremum.

  \subsubsection*{Low-Sector Spectrum Synthesis}
    \label{sec:leptonic-synthesis}

    The low-sector mass hierarchy combines two distinct mechanisms:

    \begin{itemize}
      \item a topologically determined spectral invariant
      \[
        \frac{\lambda_2}{\lambda_1}=\frac{8}{3},
      \]

      \item a dynamical stability selection rule governed by
      torsional non-commutativity.
    \end{itemize}

    For the $r=2$ sector this leads to the parameter-free relation
    \begin{equation}
      \boxed{
        \frac{m_{(2)}}{m_{(1)}}
        = \sqrt{\frac{\lambda_{2}}{\lambda_{1}}}
        \cdot \frac{3}{2\alpha}
        \cdot \frac{1}{\varphi}
      }
      \qquad\text{with}\qquad
      \frac{\lambda_2}{\lambda_1}=\frac{8}{3}.
      \label{eq:muon-ratio-final}
    \end{equation}

    The robustness of $\lambda_2/\lambda_1=8/3$ is demonstrated
    independently in Appendix~\ref{sec:spectral_ratio_derivation}.

    \paragraph*{Outlook.}
      The same commutator-driven torsional mechanism extends naturally
      to higher torsion sectors ($r\ge 3$), providing a route to the
      mass scales of composite sectors.
